% =============================================================================
%  STA6938 Time Series -- Final Project Report
%  Monthly Atmospheric CO2 at Mauna Loa, 1958--2026
%  Compile: pdflatex main && bibtex main && pdflatex main && pdflatex main
%  (On Overleaf: upload the zip and press Recompile.)
% =============================================================================
\documentclass[11pt,letterpaper]{article}

% ---- Fonts and typography ----------------------------------------------------
\usepackage[T1]{fontenc}
\usepackage[utf8]{inputenc}
\usepackage[sc]{mathpazo}          % Palatino text and math
\renewcommand{\ttdefault}{lmtt}   % scalable typewriter font
\linespread{1.04}
\usepackage{microtype}
\usepackage[margin=1in]{geometry}
\usepackage{amsmath,amssymb}
\usepackage{enumitem}
\setlist{nosep, leftmargin=*, topsep=2pt}
\setlength{\parskip}{3pt}

% ---- Colour, links, citations ------------------------------------------------
\usepackage[dvipsnames]{xcolor}
\definecolor{accent}{HTML}{1C5CAB}
\usepackage[colorlinks=true, linkcolor=accent, citecolor=accent, urlcolor=accent]{hyperref}
\usepackage[round, authoryear]{natbib}
\setlength{\bibsep}{0pt}
\renewcommand{\bibfont}{\footnotesize}

% ---- Section headings --------------------------------------------------------
\usepackage{titlesec}
\titleformat{\section}{\color{accent}\large\bfseries}{\thesection.}{0.5em}{}
\titlespacing*{\section}{0pt}{8pt}{3pt}
\titleformat{\subsection}[runin]{\bfseries}{}{0pt}{}[.]
\titlespacing*{\subsection}{0pt}{4pt}{0.6em}

% ---- Header and footer -------------------------------------------------------
\usepackage{fancyhdr}
\pagestyle{fancy}
\fancyhf{}
\fancyhead[L]{\small\textsc{sta6938 time series}}
\fancyhead[R]{\small\textsc{atmospheric co$_2$ at mauna loa}}
\fancyfoot[C]{\small\thepage}
\renewcommand{\headrulewidth}{0.4pt}

% ---- Notation ----------------------------------------------------------------
\newcommand{\co}{CO$_2$}
\newcommand{\ppm}{\,ppm}
\newcommand{\nab}{\nabla}

\begin{document}

% ---- Title block ---------------------------------------------------------------
\thispagestyle{fancy}
\begin{center}
  {\LARGE\bfseries\color{accent} Modeling and Forecasting Monthly Atmospheric \co\\[2pt] at Mauna Loa, 1958--2026}\\[8pt]
  {\large Final Project Report}\\[6pt]
  {\normalsize Student Name \quad$\cdot$\quad STA6938 Time Series \quad$\cdot$\quad Instructor: Y.~Jia \quad$\cdot$\quad December 6, 2026}
\end{center}
\vspace{-6pt}
\noindent\rule{\linewidth}{0.4pt}
\vspace{-8pt}
\begin{quote}\small
\noindent\textbf{Abstract.}
This report analyzes 822 monthly mean concentrations of atmospheric carbon dioxide recorded at Mauna Loa Observatory from March 1958 to August 2026. The series has a strong, accelerating trend and a stable 12-month cycle, so both a regular and a seasonal difference are required for stationarity. Identification from the sample ACF and PACF, together with information criteria, selects the multiplicative seasonal model ARIMA$(0,1,1)(0,1,1)_{12}$. Both moving-average parameters are highly significant, the model is invertible, and its residuals pass the Ljung--Box test at lags 6 through 24 and show constant variance. A mild departure from normality is traced to five isolated months and does not affect the estimates. In a 36-month hold-out comparison the model matches Holt--Winters smoothing and clearly beats a regression on trend and seasonal dummies. The model forecasts that monthly \co\ will first exceed 440\ppm\ in April 2030.

\noindent\textbf{Keywords:} seasonal ARIMA; differencing; Holt--Winters; Ljung--Box test; Mauna Loa; carbon dioxide.
\end{quote}
\vspace{-10pt}
\noindent\rule{\linewidth}{0.4pt}
\vspace{-14pt}

% =============================================================================
\section{Topic and Motivation}
Carbon dioxide is the most important long-lived greenhouse gas, and the Mauna Loa record, started by Charles D.~Keeling in 1958, is the longest continuous direct measurement of it \citep{keeling1960}. It is a benchmark for climate science and policy, which are framed around how fast \co\ is rising and when it will cross particular thresholds. I chose this series because its trend, seasonal cycle and short-term dependence allow nearly every method in the course to be applied to real data with a clear scientific meaning. It also raises questions that time series methods can answer quantitatively: how quickly is \co\ increasing, is the increase accelerating, and what concentrations should be expected over the next five years?

% =============================================================================
\section{Description of the Series}
The data are the monthly mean \co\ mole fractions, in parts per million (ppm), published by the NOAA Global Monitoring Laboratory \citep{noaa2026}. The series $X_t$ contains $n = 822$ consecutive months from March 1958 to August 2026, with no missing values. Concentration rose from 315.71\ppm\ in March 1958 to 427.55\ppm\ in August 2026, an increase of 111.8\ppm\ (35\%). The sample mean is 361.36\ppm\ (standard deviation 33.46\ppm). The minimum, 312.42\ppm, occurred in October 1958 and the maximum, 432.34\ppm, in May 2026. The average 12-month increase over the whole record is 1.68\ppm\ per year, but this rate has more than tripled, from 0.86\ppm/yr in the 1960s to 2.40\ppm/yr in the 2010s and 2.54\ppm/yr in the 2020s. The largest 12-month increases on record occurred in March 2024 (4.39\ppm) and April 2016 (4.25\ppm), both during strong El~Ni\~no events.

A classical additive decomposition with a $2\times12$ moving-average filter \citep[Note~4]{jia2026} attributes 96.6\% of the variance of $X_t$ to the trend. The seasonal component peaks in May (+3.04\ppm) and reaches its minimum in October ($-$3.25\ppm), an amplitude of 6.29\ppm. This cycle reflects uptake of \co\ by Northern Hemisphere vegetation in the growing season and its release in autumn and winter. The irregular component is small (standard deviation 0.30\ppm).

% =============================================================================
\section{Stationarity of the Series}
\subsection{Transformation}
A log or Box--Cox transformation is warranted when the variability of a series grows with its level \citep[Note~4]{jia2026}. Here, the within-year range of the detrended series varies only between 5.3 and 7.6\ppm. It averaged 6.0\ppm\ in the 1960s and 7.0\ppm\ in the 2010s, a 16\% increase while the level rose by 35\%. Because the seasonal swings do not grow in proportion to the level, the additive form is adequate and the series is analyzed without transformation.

\subsection{Plot and sample ACF}
The time plot shows a persistent upward trend, so the mean of $X_t$ depends on $t$ and the series is clearly non-stationary \citep[Notes~3--4]{jia2026}. The sample ACF confirms this: $\hat\rho(1) = 0.996$, and the ACF is still 0.865 at lag 36. After a first difference, $\nab X_t$, the ACF at the seasonal lags 12, 24 and 36 remains 0.925, 0.910 and 0.894. This failure to decay shows that the seasonal pattern is also non-stationary, so a seasonal difference $\nab_{12} = 1-B^{12}$ is required. The seasonally differenced series $\nab_{12}X_t$ is the 12-month change in concentration. Its ACF still decays slowly (0.856, 0.792 and 0.745 at lags 1--3), and its level rises steadily over time as the growth rate accelerates, so it does not have a constant mean.

\subsection{Augmented Dickey--Fuller tests}
Applied directly to a seasonal series at 0--2 lags, the ADF test is misleading because the seasonal swings mimic mean reversion, so it was applied to the seasonally adjusted series \citep{dickey1979}. With a single mean, $\tau$ ranges from 3.57 to 4.85, far above the 5\% critical value of $-2.86$; with a trend, $\tau$ ranges from $-0.88$ to $-0.62$, against a critical value of $-3.41$. Adding 13 augmentation lags to the raw series gives $\tau=-0.33$. In every version a unit root cannot be rejected. For $\nab_{12}X_t$ the test does reject ($\tau$ from $-7.91$ to $-5.37$), but a model for $\nab_{12}X_t$ with an AR(1) term and drift estimates $\hat\phi = 0.9993$ (SE 0.0010), practically a unit root, in line with the slowly decaying ACF and accelerating growth. For the doubly differenced series $\nab\nab_{12}X_t$, the test rejects decisively at every lag ($\tau$ from $-37.89$ to $-22.05$). I therefore model $W_t = (1-B)(1-B^{12})X_t$, which has 809 observations, as a stationary series.

\subsection{Estimation versus differencing}
As the alternative approach from Note~4, I estimated the deterministic components by regression on a quadratic trend and 11 monthly dummy variables. The fit is excellent ($R^2 = 0.9994$, residual SD 0.82\ppm), and the positive quadratic coefficient confirms that the trend is accelerating. The residuals, however, are far from white noise: their lag-1 autocorrelation is 0.915 and the Ljung--Box statistic at lag 12 is $Q(12) = 5991$. Deviations from a fixed polynomial persist for years, so the OLS standard errors and prediction intervals are unreliable (Section~7). Differencing, which lets the trend and seasonal pattern evolve stochastically, is more appropriate.

% =============================================================================
\section{Model Selection and Model-Building Procedure}
\subsection{White-noise test}
Before building a model, I checked whether $W_t$ is already white noise \citep[Note~6]{jia2026}. The Ljung--Box test \citep{ljung1978} gives $Q(6)=69.99$, $Q(12)=283.32$, $Q(18)=295.43$ and $Q(24)=301.19$, all with $p<0.001$. The differenced series therefore still has autocorrelation that a model must describe.

\subsection{Identification}
The sample ACF of $W_t$ has a significant spike at lag 1 ($-0.280$) and a large spike at lag 12 ($-0.471$), with smaller satellite values at lags 11 (0.185) and 13 (0.042). Almost all other autocorrelations lie within the $\pm1.96/\sqrt{809} = \pm0.069$ bounds. The sample PACF, by contrast, tails off: at lags 1--3 it is $-0.280$, $-0.146$ and $-0.135$, and at the seasonal lags 12, 24 and 36 it is $-0.378$, $-0.260$ and $-0.163$. An ACF that cuts off together with a PACF that tails off is the signature of moving-average behavior \citep[Notes~7--8]{jia2026}. Here that behavior appears at both the non-seasonal lag 1 and the seasonal lag 12. This points to the multiplicative model of \citet{box2015}:
\begin{equation}
(1-B)(1-B^{12})X_t = (1-\theta_1 B)(1-\Theta_1 B^{12})Z_t, \qquad Z_t \sim \mathrm{WN}(0,\sigma_Z^2),
\end{equation}
that is, ARIMA$(0,1,1)(0,1,1)_{12}$. Its right-hand side is an MA(13) with only two free parameters, which explains the small autocorrelations at lags 11 and 13.

\subsection{Estimation and selection}
Nine candidate models were estimated by maximum likelihood (ML) and compared with AIC and SBC \citep{schwarz1978}. They included models with extra AR or MA terms and with seasonal autoregressive terms. ARIMA$(0,1,1)(0,1,1)_{12}$ attains the lowest SBC (465.39). The richer ARIMA$(1,1,1)(0,1,1)_{12}$ has a slightly lower AIC (447.58 versus 451.30) but a higher SBC (466.37), and ARIMA$(0,1,2)(0,1,1)_{12}$ has SBC 467.38. Replacing or supplementing the seasonal MA term with a seasonal AR term raises the SBC by at least 6.7. A purely seasonal-AR specification, ARIMA$(0,1,1)(1,1,0)_{12}$, is clearly inadequate (SBC 683.43, and its residuals fail the Ljung--Box test). Because SBC penalizes extra parameters more heavily and the sample is large, the two-parameter model was selected. Its implied autocorrelations, $\rho(1)=-0.338$, $\rho(12)=-0.495$ and $\rho(11)=\rho(13)=0.167$, closely match the sample values above, which further supports the identification.

% =============================================================================
\section{Hypothesis Tests and Confidence Intervals}
The ML estimates are $\hat\theta_1 = 0.3885$ (SE 0.0367) and $\hat\Theta_1 = 0.8700$ (SE 0.0170), with $\hat\sigma_Z^2 = 0.1012$, so $\hat\sigma_Z = 0.318$\ppm. The log-likelihood is $-222.65$. For $H_0:\theta_1=0$, $z = 10.59$ ($p<0.001$), with 95\% confidence interval $[0.317,\,0.460]$. For $H_0:\Theta_1=0$, $z = 51.28$ ($p<0.001$), with 95\% interval $[0.837,\,0.903]$. Both terms are therefore essential. Conditional least squares \citep[Notes~7--9]{jia2026} gives similar values, $\hat\theta_1 = 0.4172$ (SE 0.0360) and $\hat\Theta_1 = 0.8116$ (SE 0.0194). This shows that the conclusions do not depend on the estimation method.

The root of $1-\hat\theta_1B$ is $1/\hat\theta_1 = 2.57$, and the roots of $1-\hat\Theta_1B^{12}$ have modulus $\hat\Theta_1^{-1/12} = 1.012$. Both lie outside the unit circle, so the fitted model is \emph{invertible} and can be used for forecasting. The non-seasonal $\pi$ weights ($0.389$, $0.151$, $0.059$) decay quickly. The confidence interval for $\Theta_1$ lies entirely below 1. A value of $\Theta_1$ equal to 1 would mean the seasonal difference had been applied to a fixed, deterministic seasonal pattern, so the interval shows that seasonal differencing was justified. The large value of $\hat\Theta_1$ means the seasonal pattern is very stable and changes only slowly from year to year. The smaller $\hat\theta_1$ means about 39\% of each month's shock is reversed the following month.

% =============================================================================
\section{Model Adequacy Checking and Remedial Methods}
\subsection{Independence}
The residuals pass the Ljung--Box test at the lags used in class: $Q(6) = 9.05$ ($p=0.060$), $Q(12) = 9.74$ ($p=0.464$), $Q(18)=14.43$ ($p=0.567$) and $Q(24)=20.83$ ($p=0.532$), with degrees of freedom reduced by the two estimated parameters. Only one of the first 36 residual autocorrelations lies outside $\pm0.069$: lag 3, at $-0.081$. About 1.8 such exceedances are expected by chance among 36 lags of white noise.

\subsection{Constant variance}
The Ljung--Box test on the squared residuals gives $Q(12)=12.40$ ($p=0.414$) and $Q(24)=35.36$ ($p=0.063$). The residual variance is similar in the first and second halves of the sample (0.091 and 0.107). There is no evidence that the variance changes over time.

\subsection{Normality}
The residuals are nearly symmetric, with skewness 0.11 and excess kurtosis 0.35, and the normal Q--Q plot is close to linear. However, three tests reject normality at the 5\% level: Shapiro--Wilk $W = 0.9963$ ($p=0.048$) \citep{shapiro1965}, Anderson--Darling ($p=0.038$) and Jarque--Bera ($p=0.006$). Five months have standardized residuals beyond $\pm3$: March 1959 ($-3.46$), March 1977 (3.31), September 1986 (3.11), April 2016 (3.97) and March 2022 ($-3.53$). The April 2016 value coincides with the unusually rapid \co\ growth during the 2015--16 El~Ni\~no.

\subsection{Remedial methods}
\begin{enumerate}[label=(\roman*)]
  \item \emph{Lag-3 autocorrelation.} Following the model-refinement step of Note~9, I added an AR(1) term. ARIMA$(1,1,1)(0,1,1)_{12}$ has a significant $\hat\phi_1 = 0.215$ (SE 0.084, $p=0.010$) and reduces the lag-3 residual autocorrelation to $-0.042$. Its AR and MA roots (4.65 and 1.74) are distinct, so there is no parameter redundancy. However, it raises the SBC, and in the hold-out test of Section~7 its forecasts were less accurate (RMSE 1.168 versus 1.110) and its 95\% intervals covered only 83\% of outcomes, against 97\% for the simpler model. The extra parameter fits noise rather than structure, so the parsimonious model was retained.
  \item \emph{Non-normality.} The five outlying months were treated as missing and the model was re-estimated by exact likelihood (Kalman filter). The estimates were almost unchanged ($\hat\theta_1 = 0.3838$, $\hat\Theta_1 = 0.8689$), and normality was no longer rejected (Shapiro--Wilk $p=0.211$, Anderson--Darling $p=0.101$, Jarque--Bera $p=0.094$). The residuals still passed the Ljung--Box tests ($p \ge 0.092$). The non-normality is thus caused by a handful of genuine but unusual months. With $n=822$ the ML estimates are approximately normal by large-sample theory, so the tests and intervals above remain valid, and the full-data model is kept because there is no reason to discard real observations.
  \item \emph{Transformation.} As discussed in Section~3, the residual variance is stable over time, so no variance-stabilizing transformation is needed.
\end{enumerate}
The fitted ARIMA$(0,1,1)(0,1,1)_{12}$ model is therefore adequate.

% =============================================================================
\section{Forecasting}
\subsection{Out-of-sample comparison}
To assess predictive performance, I re-estimated each method on data up to August 2023 and forecast the remaining 36 months (September 2023 to August 2026). The comparison includes the selected model, additive Holt--Winters exponential smoothing \citep{winters1960}, the regression on trend and seasonal dummies, and the seasonal naive method. The selected model achieved RMSE 1.110\ppm\ and MAPE 0.24\%, and its 95\% prediction intervals covered 97.2\% of the observed values. Holt--Winters, the seasonal method of Note~5, performed similarly (RMSE 0.981\ppm, coverage 97.2\%). The regression forecasts were less accurate (RMSE 1.491\ppm), and its intervals covered only 52.8\% of outcomes because they ignore the strong residual autocorrelation. The seasonal naive method was much worse (RMSE 6.448\ppm). All methods under-predicted 2024, when the annual mean rose by a record 3.52\ppm. This is consistent with the 2023--24 El~Ni\~no, an external effect that no univariate model can anticipate. On the full data, Holt--Winters estimates $\hat\alpha = 0.534$ (level), $\hat\beta = 0.014$ (trend) and $\hat\gamma = 0.299$ (seasonal), with a current slope of 0.203\ppm\ per month (about 2.4\ppm\ per year), and its one-step in-sample RMSE (0.324\ppm) is close to the ARIMA model's (0.315\ppm).

\subsection{Forecasts to 2030}
The model estimated on all 822 months gives a forecast for September 2026 of 426.16\ppm, with 95\% prediction interval [425.53, 426.78]. For the May 2027 seasonal peak it gives 434.36\ppm\ [433.11, 435.60]. The forecast first exceeds 440\ppm\ in April 2030 (440.84\ppm\ [437.78, 443.89]) and peaks at 441.59\ppm\ [438.49, 444.69] in May 2030. The forecast for December 2030 is 439.03\ppm\ [435.59, 442.47]. The forecast annual means, 431.44 (2027), 433.85 (2028), 436.26 (2029) and 438.67\ppm\ (2030), imply continued growth of about 2.4\ppm\ per year. The 95\% interval widens from 1.25\ppm\ one month ahead to 6.88\ppm\ at the end of 2030, reflecting the accumulated uncertainty of a doubly integrated process. The independently estimated Holt--Winters forecast for December 2030 (439.16\ppm) agrees with the ARIMA forecast to within 0.13\ppm.

% =============================================================================
\section{Conclusions}
Monthly \co\ at Mauna Loa is a non-stationary series with an accelerating stochastic trend and a stable additive seasonal cycle. Regular and seasonal differencing produce a stationary series, and the ACF/PACF, the information criteria and the residual diagnostics all support the parsimonious ARIMA$(0,1,1)(0,1,1)_{12}$ model. The fitted model is invertible, its two parameters are estimated precisely, its residuals are uncorrelated with constant variance, and its inferences are robust to the few outlying months. The data show that \co\ is not only increasing but increasing faster: the growth rate rose from under 1\ppm\ per year in the 1960s to about 2.5\ppm\ per year today. The model projects that concentrations will cross 440\ppm\ in spring 2030. Two limitations remain. The model is univariate, so it cannot anticipate El~Ni\~no-driven accelerations such as those of 2016 and 2024, and it assumes that the current emissions-driven growth continues. Adding climate indices or emissions data as explanatory variables is a natural extension.

\vspace{1pt}
\noindent{\small\textit{Supplementary material.} Figures, tables and R code (R~4.3.3; \texttt{forecast} \citep{hyndman2008}, \texttt{tseries}, \texttt{urca}, \texttt{nortest}) reproducing every number are provided separately.}

\bibliographystyle{plainnat}
{\small\bibliography{references}}

\end{document}
